\documentclass[12pt,letterpaper]{report}
\usepackage{graphicx}
\usepackage{scrextend}
\usepackage{vmargin}
\usepackage{graphicx}
\usepackage{multirow}
\usepackage[utf8]{inputenc}
\usepackage[spanish]{babel}
\usepackage{multicol}
\usepackage{enumerate}
\usepackage{float}
\usepackage{amsmath, amsthm, amssymb, amsfonts}
\usepackage[usenames]{color}
\parindent=0mm
\pagestyle{empty}
\definecolor{miorange}{rgb}{0.91, 0.43, 0.0}
\begin{document}
\setmargins{2.5cm}      
{1.5cm}                     
{2cm}  
{24cm}                    
{10pt}                          
{1cm}                          
{0pt}                             
{2cm}
\begin{flushright}
\today
\end{flushright}
\vspace{-1.75cm}
\section*{Tarea Clase 14}
\subsection*{Obtenga las soluciones de las siguientes ecuaciones:}
\begin{enumerate}
\item Un chef va a preparar una ensalada de verduras con tomate, zanahoria, papa y brócoli. ¿De cuántas formas se puede preparar la ensalada usando solo 2 ingredientes?
\item Eduardo, Carlos y Sergio se han presentado a un concurso de pintura. El concurso otorga \$200 al primer lugar y  \$100 al segundo. ¿De cuántas formas se pueden repartir los premios de primer y segundo lugar?
\item Se va a programar un torneo de ajedrez para los 10 integrantes de un club. ¿Cuántos partidos se deben programar si cada integrante jugará con cada uno de los demás sin partidos de revancha?
\item Se va a programar un torneo de ajedrez para los 10 integrantes de un club. ¿Cuántos partidos se deben programar si cada integrante jugará con cada uno de los demás sin partidos de revancha?
\item ¿De cuántas formas distintas pueden sentarse ocho personas alrededor de una mesa redonda?
\item ¿Cuántos números de cinco cifras distintas se pueden formar con las cifras impares? ¿Cuántos de ellos son mayores de 70.000?
\item Una mesa presidencial está formada por ocho personas ¿De cuántas formas distintas se pueden sentar, si el presidente y el secretario siempre van juntos?
\end{enumerate}
Las respuestas de las tareas pueden ser enviadas al siguiente correo giovannilopez9808@gmail.com con el asunto entre corchetes de la siguiente manera [Tarea \#] y el nombre del archivo con sus apellidos y nombre, ya sea el nombre completo o solo una parte, por ejemplo \textit{LopezPadilla.pdf}, \textit{LopezGiovanni.pdf}, \textit{LopezPadillaGiovanniGamaliel.pdf}. Las respuestas envienlas en formato pdf, estas pueden ser totalmente digitales (que escriban todo en la computadora) o que sean fotografías, pero por favor, envielo en solo un documento.
\end{document}